\section{Collisionless matter in a dynamical spacetime}\label{sec:system}

Stars in dense clusters, particles in dark-matter spikes, and debris in colliding stellar streams move freely through smooth collective gravitational fields \cite{1985ApJ...298...34S,2008gady.book.....B,Andreasson:2011ng}. Two-body gravitational encounters act over timescales far longer than a dynamical crossing time, so each constituent follows a geodesic of the spacetime geometry generated by the whole ensemble. Throughout this paper we set Newton's gravitational constant $G$ and the speed of light $c$ to unity and adopt metric signature $(-,+,+,+)$. Greek letters $\mu,\nu,\rho,\sigma$ label spacetime indices, Latin letters $i,j,k,l$ label spatial coordinate indices, and parenthesized indices $(m),(n),(k)$ label components in a tetrad, a local frame of four mutually perpendicular unit vectors.

Collisionless matter in a curved spacetime with metric $g_{\mu\nu}$ and inverse $g^{\mu\nu}$ is described in phase space by a rest-mass-weighted distribution function $f(x^\mu, u^\nu)$ \cite{Andreasson:2011ng,Sarbach:2013fya,Acuna-Cardenas:2021nkj}. Here $x^\mu$ are spacetime coordinates and $u^\mu$ is the four-velocity, equal to the four-momentum per unit rest mass on the future mass shell $g_{\mu\nu} u^\mu u^\nu = -1$. Phase-space conservation along free-fall worldlines coupled to Einstein's field equations forms the Einstein-Vlasov system,
\begin{align}\label{eq:einstein_vlasov}
u^\mu \partial_\mu f - \Gamma^\mu_{\nu\rho} u^\nu u^\rho \frac{\partial f}{\partial u^\mu} &= 0, \nonumber\\
G_{\mu\nu} &= 8\pi T_{\mu\nu},
\end{align}
where $\Gamma^\mu_{\nu\rho}$ are the Christoffel symbols of $g_{\mu\nu}$. The Einstein tensor $G_{\mu\nu}$ is sourced by the stress-energy tensor $T_{\mu\nu}$, given by the second four-velocity moment of $f$ over the future mass shell.

Time evolution splits the spacetime into spacelike hypersurfaces of constant coordinate time $t$ with spatial coordinates $x^i$, induced metric $\gamma_{ij}$, inverse $\gamma^{ij}$, and determinant $\gamma = \det(\gamma_{ij})$ \cite{Arnowitt:1962hi}. Local measurements on each slice are referred to normal observers, Eulerian observers whose worldlines cross the slice orthogonally with future-directed unit four-velocity $n^\mu$ and covector $n_\mu = (-\alpha, 0, 0, 0)$. The lapse $\alpha$ links adjacent slices by converting a coordinate interval $\dd t$ into proper time $\alpha\,\dd t$ along normal observers. The shift $\beta^i$ gives the coordinate velocity at which spatial grid points drift relative to those normal observers. In these variables the spacetime line element reads
\begin{equation}\label{eq:line_element}
\dd s^2 = -\alpha^2\,\dd t^2 + \gamma_{ij}(\dd x^i + \beta^i\,\dd t)(\dd x^j + \beta^j\,\dd t),
\end{equation}
with $n^\mu = (\alpha^{-1}, -\alpha^{-1}\beta^i)$. The bending of each slice inside the ambient spacetime is measured by the symmetric extrinsic curvature $K_{ij}$, defined through
\begin{equation}\label{eq:extrinsic_curvature}
\partial_t \gamma_{ij} = -2\alpha K_{ij} + D_i \beta_j + D_j \beta_i,
\end{equation}
where $D_i$ is the spatial covariant derivative compatible with $\gamma_{ij}$ and $\beta_i = \gamma_{ij}\beta^j$. Normal observers project $T^{\mu\nu}$ into ten matter fields---the energy density $E$, covariant momentum density $S_i$, and symmetric spatial stress $S_{ij}$---according to
\begin{align}\label{eq:source_projections}
E &= n_\mu n_\nu T^{\mu\nu}, &
S_i &= -g_{i\mu} n_\nu T^{\mu\nu}, \nonumber\\
S_{ij} &= g_{i\mu} g_{j\nu} T^{\mu\nu}. & &
\end{align}

Direct grid discretization of $f$ spans six phase-space dimensions and becomes prohibitive in three spatial dimensions without symmetry, whereas the Vlasov equation conserves phase-space weight along timelike geodesics. A particle-in-cell method \cite{1988csup.book.....H,1991ppcs.book.....B,Wu:2026nrpic} exploits these characteristics by sampling $f$ with $N$ finite-weight phase-space packets called macroparticles, indexed by $p = 1, \dots, N$. With $u_\mu = g_{\mu\nu} u^\nu$, each macroparticle carries fixed rest mass $m_p$, spatial coordinates $x_p^i(t)$, and covariant spatial four-velocity components $u_{i,p}(t) = g_{i\mu} u_p^\mu$, giving six phase-space variables per particle. On the future mass shell $g^{\mu\nu} u_\mu u_\nu = -1$, the Lorentz factor seen by normal observers is
\begin{equation}\label{eq:lorentz_factor}
W = \alpha u^0 = \sqrt{1 + \gamma^{ij} u_i u_j},
\end{equation}
which determines the single-particle Hamiltonian $-u_0 = \alpha W - \beta^i u_i$. Hamilton's canonical equations for $(x^i, u_i)$ yield the characteristic equations,
\begin{align}\label{eq:characteristics}
\frac{\dd x^i}{\dd t} &= \frac{\alpha \gamma^{ij} u_j}{W} - \beta^i, \nonumber\\
\frac{\dd u_i}{\dd t} &= -W \partial_i \alpha + u_j \partial_i \beta^j - \frac{\alpha}{2W} u_j u_k \partial_i \gamma^{jk}.
\end{align}
At spatial coordinate $\bm{x} = (x^1, x^2, x^3)$, the ensemble stress-energy projected onto normal observers forms the particle sources,
\begin{align}\label{eq:particle_sources}
E(t,\bm{x}) &= \sum_{p=1}^{N} \frac{m_p W_p}{\sqrt{\gamma}}\,\delta^3(\bm{x} - \bm{x}_p), \nonumber\\
S_i(t,\bm{x}) &= \sum_{p=1}^{N} \frac{m_p u_{i,p}}{\sqrt{\gamma}}\,\delta^3(\bm{x} - \bm{x}_p), \nonumber\\
S_{ij}(t,\bm{x}) &= \sum_{p=1}^{N} \frac{m_p u_{i,p} u_{j,p}}{W_p \sqrt{\gamma}}\,\delta^3(\bm{x} - \bm{x}_p).
\end{align}
Here $\delta^3(\bm{x} - \bm{x}_p)$ is the three-dimensional coordinate Dirac distribution and $W_p$ is evaluated at $\bm{x}_p(t)$. Geodesic transport of $(x_p^i, u_{i,p})$ under the characteristic equations, coupled to smoothed mesh projections of the particle sources, solves the Einstein-Vlasov system.
